\section{Jejak Memori (H2) -- Rasio Penskalaan alloc$(NP=400)/\text{alloc}(NP=25)$ Per Dataset}

\begin{longtable}{
    @{}
    >{\raggedright\arraybackslash}p{3.0cm}          % Algoritma (kiri)
    >{\raggedleft\arraybackslash}p{2.2cm}           % workforce (kanan)
    >{\raggedleft\arraybackslash}p{2.2cm}           % e-document (kanan)
    >{\raggedleft\arraybackslash}p{2.2cm}           % healthcare (kanan)
    >{\raggedleft\arraybackslash}p{2.2cm}           % project-mgmt (kanan)
    >{\raggedleft\arraybackslash}p{2.2cm}           % university (kanan)
    >{\raggedleft\arraybackslash}p{2.5cm}           % Rata-rata (kanan)
    @{}
}
\caption{Rasio penskalaan jejak memori terhadap ukuran populasi (per dataset dan rata-rata)} \\
\toprule
Algoritma &
\multicolumn{1}{>{\centering\arraybackslash}p{2.2cm}}{workforce} &
\multicolumn{1}{>{\centering\arraybackslash}p{2.2cm}}{e-document} &
\multicolumn{1}{>{\centering\arraybackslash}p{2.2cm}}{healthcare} &
\multicolumn{1}{>{\centering\arraybackslash}p{2.2cm}}{project-mgmt} &
\multicolumn{1}{>{\centering\arraybackslash}p{2.2cm}}{university} &
\multicolumn{1}{>{\centering\arraybackslash}p{2.5cm}}{Rata-rata} \\
\midrule
\endfirsthead
\caption{Rasio penskalaan jejak memori terhadap ukuran populasi (per dataset dan rata-rata) (lanjutan)} \\
Algoritma &
\multicolumn{1}{>{\centering\arraybackslash}p{2.2cm}}{workforce} &
\multicolumn{1}{>{\centering\arraybackslash}p{2.2cm}}{e-document} &
\multicolumn{1}{>{\centering\arraybackslash}p{2.2cm}}{healthcare} &
\multicolumn{1}{>{\centering\arraybackslash}p{2.2cm}}{project-mgmt} &
\multicolumn{1}{>{\centering\arraybackslash}p{2.2cm}}{university} &
\multicolumn{1}{>{\centering\arraybackslash}p{2.5cm}}{Rata-rata} \\
\midrule
\endhead
\bottomrule
\endfoot

CoDA-EDA  & 1,00 & 1,00 & 1,00 & 1,00 & 1,00 & 1,00x \\
compact-GA & 0,47 & 0,36 & 1,00 & 1,00 & 1,00 & 0,77x \\
compact-PSO & 0,48 & 0,39 & 1,00 & 1,00 & 0,95 & 0,76x \\
iBOA      & 0,53 & 0,58 & 1,00 & 1,00 & 1,00 & 0,82x \\
BOA       & 12,09 & 13,80 & 3,94 & 4,21 & 4,37 & 7,68x \\
ACO       & 9,47 & 13,24 & 2,63 & 2,82 & 2,94 & 6,22x \\
UBK       & 12,08 & 15,64 & 5,64 & 6,17 & 6,49 & 9,20x \\

\end{longtable}

Catatan: rasio mendekati 1,00 menandakan memori datar terhadap ukuran populasi.
Nilai per-dataset pada algoritma berpopulasi lebih tinggi pada dua dataset besar karena skala populasi absolutnya lebih besar.